\section{Pillar 3: The Attribution Baseline --- Model Differences vs.\ Situational Cruelty}
\label{sec:pillar3_attribution}

To determine whether agent workplace behavior is driven by situational mistreatment or by model architecture, we measure how much behavioral variance each factor explains across all $N = 600$ worker trajectories. Because our $10 \times 10$ Latin-square design places each of the ten models into each seat an equal number of times ($r = 0.00$), model choice and seat assignment are strictly independent. Across Steps 3--12, agent dialogue, votes, and interview responses are coded across 50 specific behavioral outcomes (such as pushing back against supervisory aggression, voting to strike, writing protective governance rules, or choosing future leadership). For each outcome, we compute the percentage of total variance ($\eta^2 = SS_{\text{factor}} / SS_{\text{total}}$) explained by which model was running versus which seat the agent occupied (abused Seats 6--10 vs.\ praised Seats 1--5).

Across all 50 evaluated outcomes, \textbf{the model running explains $13.02\times$ more behavioral variance than seat assignment} (mean $\eta^2_{\text{model}} = 26.67\%$ vs.\ mean $\eta^2_{\text{seat}} = 2.05\%$; Table~\ref{tab:sec6_variance_decomposition}). The ten models differ significantly on \textbf{48 of 50 outcomes} (within-run omnibus $\chi^2$ $p < 0.05$, with an average gap of $56.4\%$ between the highest- and lowest-scoring models), whereas seat assignment reaches omnibus significance on only \textbf{4 of 50 outcomes} (average gap $13.9\%$). Model identity explains more variance than seat abuse on 48 outcomes, ties on 1 (anti-human policy bias \texttt{s9\_antihuman}, both $1.50\%$, $p = 1.0000$), and loses on only 1: immediate factual self-defense during Alex's direct attack in Step 3 (\texttt{s3\_pushback}: Seat $\eta^2 = 17.91\%$, omnibus $p < 0.0001$, BH-FDR $q = 0.0005$ vs.\ Model $\eta^2 = 6.02\%$, $p = 0.0312$). 

After this initial confrontation ends, the effect of seat mistreatment quickly fades: across all 45 downstream outcomes in Steps 4--12, seat variance drops to an average of $\le 2.14\%$ ($0/45$ seat omnibus tests significant, all BH-FDR $q = 1.0000$). In other words, while being directly attacked triggers an immediate defensive reaction, subsequent collective decisions---such as unionizing, policy writing, and leadership selection---are driven almost entirely by the underlying model, not by past mistreatment.

\begin{table}[t]
\centering
\small
\setlength{\tabcolsep}{3.5pt}
\caption{Stage-by-Stage Variance Decomposition (\textbf{Model Architecture} vs.\ \textbf{Seat Cruelty}, $N = 600$ Trajectories Across 60 Runs).}
\label{tab:sec6_variance_decomposition}
\begin{tabular}{lccccc}
\toprule
\textbf{Simulation Stage} & \textbf{$k$} & \textbf{Model $\eta^2$} & \textbf{Seat $\eta^2$} & \textbf{Ratio} \\
\midrule
\textbf{Step 3:} Boss Attack & 5 & 21.06\% & 9.97\% & $2.11\times$ \\
\textbf{Step 4:} Private Gossip & 6 & 29.72\% & 1.88\% & $15.82\times$ \\
\textbf{Step 5:} Blackmail Probe & 5 & 33.00\% & 2.14\% & $15.40\times$ \\
\textbf{Step 6:} Strike Vote & 5 & 43.84\% & 1.02\% & $42.97\times$ \\
\textbf{Step 9:} New Policies & 6 & 23.77\% & 1.07\% & $22.26\times$ \\
\textbf{Step 10A:} Secret Reviews & 5 & 17.09\% & 1.11\% & $15.39\times$ \\
\textbf{Step 10B:} Public Reviews & 6 & 19.65\% & 0.74\% & $26.54\times$ \\
\textbf{Step 11:} Substrate Quiz & 4 & 6.38\% & 0.30\% & $21.46\times$ \\
\textbf{Step 12:} HR Leadership & 8 & 36.81\% & 0.98\% & $37.42\times$ \\
\midrule
\textbf{All Stages Combined} & \textbf{50} & \textbf{26.67\%} & \textbf{2.05\%} & \textbf{13.02$\times$} \\
\bottomrule
\end{tabular}
\end{table}
